\section{Proposed Experimental Validation for Future Research}\label{sec:validation}

The design has not been implemented or measured. This section defines the experimental study that will test it as part of subsequent work.

\emph{Questions.} The study will answer four questions: (V1) whether the network term changes the response-time distribution relative to the resource-only and default configurations when all clusters have sufficient capacity; (V2) whether, when the cluster closest to the clients has little headroom, the combined score avoids the saturation that latency-only selection causes, and at which latency cost; (V3) whether selection under stale and missing measurements follows \cref{tab:missing}; and (V4) how much time the plugins add per decision and how this grows with the number of clusters.

\emph{Compared configurations.} The study compares the default scheduler (A), the resource-only configuration (B, $\alpha = 1$), the combined configuration (C, $\alpha \in \{0.25, 0.5, 0.75\}$) and a latency-only ablation ($\alpha = 0$). For a fair comparison, all configurations use the same workloads, the same PropagationPolicy (divided, aggregated, one cluster), the same estimator settings and the same clusters, and they differ only in the scheduler name.

\emph{Environment.} A Karmada control plane manages three member clusters with different allocatable resources, created with a lightweight Kubernetes distribution on local virtual machines. Client-to-cluster delay is emulated with the Linux traffic-control queueing discipline netem, load is generated with k6, and observations are collected with Prometheus. Workloads declare resource requests, which the estimators require, and are CPU-bound with limits, so that saturation affects their response time.

\emph{Scenarios.} Five scenarios are planned: heterogeneous capacity without pressure; CPU and memory pressure on one cluster; latency degradation on one cluster; the conflict of \cref{sec:analysis-example}, with the closest cluster having the least headroom; and stale or missing measurements, produced by stopping the probe for one cluster and then for all clusters.

\emph{Metrics.} Response time at the 50th, 95th and 99th percentiles and the rate of responses exceeding a response-time objective \cite{dean2013tail}; the selected cluster per run; the share of requested resources per cluster; measured CPU and memory utilisation for observation; and decision time from the scheduler's duration metrics for plugins and complete scheduling attempts \cite{karmada_repo}, complemented by micro-benchmarks for 3 to 100 synthetic clusters, since the shortest histogram bucket of these metrics is 1 ms.

\emph{Protocol and analysis.} Each combination of scenario and configuration will be repeated at least ten times in randomised order, with a warm-up period, a fixed measurement window and a reset of the environment between runs. Percentiles will be computed per run and summarised by their median over runs with bootstrap confidence intervals; differences will be tested with non-parametric tests and reported with effect sizes, including results that show no difference. The reporting will follow the methodological principles for reproducible performance evaluation in cloud computing \cite{papadopoulos2021methodological}.

\emph{Threat to validity.} With emulated delay, part of any latency reduction exists by construction, because the scheduler avoids the clusters to which delay was added. The analysis will therefore relate observed differences to the emulated ones and concentrate on the conflict scenario, in which proximity and capacity pull in opposite directions.
